\section{Introduction}

Visual speech recognition (VSR) recovers spoken language from facial motion.
Recent systems learn strong representations from large audiovisual collections and
decode complete utterances with CTC, attention, or combinations of the two
\citep{graves2006ctc,watanabe2017hybrid,ma2023autoavsr}.
Complete-utterance recognition is an important benchmark, but an interface need
not request an entire utterance. A video editor may select a short passage for
transcription, or a viewer may request the speech in a particular time interval.
When audio is unavailable or unsuitable for the task, a visual recognizer must
decide both what context to observe and which recognized words to return.

The obvious solution is to crop the requested interval and run an existing
recognizer. This makes the input boundary coincide with the requested output
boundary. The coincidence is convenient computationally, but may discard useful
speech immediately outside the selection. Neighboring video can help resolve
the target's visual ambiguity or preserve temporal structure near the crop
boundary. However, passing a larger clip to the recognizer introduces another
problem: the resulting transcript includes speech that was not requested.
Context admission and output selection therefore need to be considered together.

We formulate \emph{context-aware interval transcription}: given a requested
interval and an allowance of neighboring video, return the transcript assigned
to the request. The formulation separates two decisions: which frames support
recognition, and where to enforce the output boundary. Contextual direct decoding
restricts contextual encoder states before decoding. Our
\emph{prediction-aligned projection} instead decodes the allowed video, aligns
its predicted tokens with native visual CTC scores, and selects complete words
after recognition. Both procedures use the same contextual input and frozen
recognizer. Their comparison reveals how the stage of output restriction
interacts with decoding.

Three research questions organize the study. First, does genuine neighboring
video improve recognition of a fixed target beyond the effect of input length?
Second, when does selecting from a contextual transcript help beyond selecting
contextual encoder states? Third, how do the context available around a request
and the reuse of earlier computation affect the accuracy--cost tradeoff?
Crop-first recognition, matched decoders, identical-transcript allocation
controls, and full-clip timestamp lookup provide the corresponding comparisons.
The full-clip alternative is particularly relevant when an application can
process the entire video or already holds a transcript index.

The study contributes an interval-transcription formulation and a concrete
prediction-aligned implementation that adapts pretrained VSR to selected-passage
requests. A controlled four-corpus evaluation then separates contextual encoding,
decoder behavior, and temporal output selection. Finally, context interventions,
duration grids, and repeated-query measurements establish an empirical account
of when bounded recognition is useful and when reusable full-clip processing
is preferable. Together, these contributions support passage-level visual
transcription through an explicit treatment of context, output scope, and work
per request.
